\documentclass[10pt,a4paper]{amsart}
\usepackage[margin=19mm]{geometry}
\usepackage{amsmath,amssymb,mathtools}
\usepackage[scr=rsfs,cal=euler]{mathalfa}
\usepackage{xfrac}
\usepackage[unicode,hidelinks,bookmarks=false]{hyperref}
\newcommand{\ie}{i.e.\ }
\newcommand{\cf}{cf.\ }
\newcommand{\resp}{respectively\ }
\newcommand{\Subsection}{\subsection}
\providecommand{\nobreakdash}{\nobreak\hbox{-}}
\setlength{\emergencystretch}{2em}
\multlinegap0pt
\usepackage{bm}
\usepackage{bbm}
\usepackage{dsfont}
\usepackage{xspace}
\usepackage{cancel}
\usepackage{mathtools}
\usepackage{amsthm}
\makeatletter
\@ifundefined{theorem}{\newtheorem{theorem}{Theorem}}{}
\@ifundefined{lemma}{\newtheorem{lemma}[theorem]{Lemma}}{}
\makeatother
\newcommand{\R}{\mathbb{R}}
\title[Quantitative Runge type approximation theorems for zero solu]{Quantitative Runge type approximation theorems for zero solutions of certain partial differential operators}
\author{Andreas Debrouwere, Thomas Kalmes}
\date{Source: arXiv:2209.10794v2 (CC BY 4.0)}
\begin{document}
\maketitle

\noindent\emph{Source and licence.} Quantitative Runge type approximation theorems for zero solutions of certain partial differential operators. Version arXiv:2209.10794v2 (CC BY 4.0), distributed under the Creative Commons Attribution 4.0 International licence, as recorded in the arXiv licence metadata for this exact version and shown on the abstract page https://arxiv.org/abs/2209.10794v2. Full rights notice: licenses/arxiv-2209.10794-CC-BY-4.0.txt.\par\smallskip

\noindent\textit{Editorial scope.} Counted unit: the Lemma whose source label is \texttt{lem: entailing minimum principle} in the article (source0.tex lines 320--322, source label \texttt{lem: entailing minimum principle}), delivered together with the article's own complete original proof of it (source0.tex lines 324--349, one \texttt{proof} environment of 26 lines). No mathematics is written, repaired or re-derived: statement and proof are the article's own text line for line. Required context retained uncounted: none. Every definition, display and label of those lines is retained unchanged. Omitted from this package and not redistributed: 1--319, 323--323, 350--984. Nothing inside a retained excerpt is omitted or altered: every retained body is delivered exactly as the article's lines are written, so no omission receipt of any kind accompanies this package. Citations: the retained excerpts carry no citation command at all, so no bibliography entry of the article is transcribed and no reference is redistributed. Reference adaptation: none was needed and none was made; every reference inside the delivered unit resolves inside it, so no unresolved reference or citation is produced. Printed slips kept literal: no printed word, symbol, hypothesis or number of the counted statement or of its proof was altered. Layout and class: the article's own class is \texttt{amsart} with packages including geometry, graphicx, latexsym, comment, tikz, mathrsfs, amsfonts, amssymb, amsmath, amsthm, amsxtra, amscd, hyperref, mathtools; the wrapper is the lane's pinned \texttt{amsart} editorial wrapper at 10pt on A4, which restates the theorem-like environments the retained text needs and adds the article's own macro definitions that the retained text needs (\texttt{\textbackslash R}), with extra wrapper packages (bm, bbm, dsfont, xspace, cancel, mathtools). The delivered numbering is the unit's own. Assets: no retained span contains a figure environment, a graphics-inclusion command, a PDF-page inclusion, a table or a TikZ picture; no image file is copied into this package, the package declares no asset dependency and no image file is redistributed with it. Licence: version arXiv:2209.10794v2 of this article is distributed under the Creative Commons Attribution 4.0 International licence, as recorded in the arXiv licence metadata for this exact version and shown on the abstract page https://arxiv.org/abs/2209.10794v2. A full rights notice is delivered at licenses/arxiv-2209.10794-CC-BY-4.0.txt. Characters: every retained excerpt is pure ASCII, so the delivered engine, XeLaTeX with its default Latin Modern text font, reports no missing character.
\begin{lemma}\label{lem: entailing minimum principle}
	Let $Y \subseteq X\subseteq\R^d$ be open. Let $F\subseteq\R^d$ be  a non-empty closed set such that $d_X$ satisfies the minimum principle in $x+F$ for every $x\in\R^d$. Suppose that there is no $x\in\R^d$ such that $X$ contains a compact connected component of $(\R^d\backslash Y)\cap (x+F)$. Then, $d_Y$ satisfies the minimum principle in $x+F$ for every $x\in\R^d$.
\end{lemma}

\begin{proof}
	We may assume $Y\neq X$. We argue by contradiction, so assume that there is $x\in \R^d$ and  $K\subseteq Y\cap(x+F)$ compact such that
$$
		\min_{z\in K}d_Y(z)<\min_{z\in\partial_{x+F}K}d_Y(z).
$$
	Choose $z_0\in K$ and $y_0\in\R^d\backslash Y$ with $|z_0-y_0|=\min_{z\in K}d_Y(x)$.
Hence,
	\begin{equation}\label{eq: violation minimum principle 2}
		|z_0-y_0|<\inf\{|y-z| \,; \, y\in\R^d\backslash Y, z\in \partial_{x+F}K\}.
	\end{equation}
	This implies that $\partial_{y_0-z_0+x+F}(y_0-z_0+K) = y_0-z_0+\partial_{x+F}K \subseteq Y$. Since $y_0 \in y_0-z_0+x+F$, we obtain that the connected component $C_{y_0}$ of $(\R^d\backslash Y)\cap (y_0-z_0+x+F)$ containing $y_0$ satisfies $C_{y_0}\subseteq y_0-z_0+K$. In particular, $C_{y_0}$ is compact. If $X=\R^d$, this gives the desired contradiction. 	If $X\neq\R^d$, we conclude that $\emptyset\neq\partial X\cap C_{y_0}$. Set
	$$t_0=\inf\{t\in[0,\infty) \, ; \,t(y_0-z_0)+K\cap\partial X\neq \emptyset\}.$$
 We have that $t_0 > 0$ because $K \subseteq X$ is compact, while  $C_{y_0}\subseteq y_0-z_0+K$ and $\emptyset\neq\partial X\cap C_{y_0}$ imply that $t_0 \leq 1$. Note that $t(y_0-z_0)+K\subseteq X$ for all $t \in (0,t_0)$. Since $t_0(y_0-z_0)+K\cap\partial X\neq \emptyset$, it holds that
	\begin{equation}\label{eq: contradiction minimum principle 1}
	 d\left(t(y_0-z_0)+K, \R^d\backslash X\right)\leq (t_0-t)|y_0-z_0|, \qquad 	\forall\,t\in (0,t_0).
	\end{equation}
	Set 
	$$\delta=\inf\{|y-z|\, ; \, y\in\R^d\backslash Y, z\in \partial_{x+F}K\}-|z_0-y_0|,$$
	so that $\delta>0$ by \eqref{eq: violation minimum principle 2}. For all $y\in \R^d\backslash X\subseteq\R^d\backslash Y$, $z\in \partial_{x+F}K$ and $t\in (0,t_0)$ we have that
	$$|y-\left(t(y_0-z_0)+z\right)|\geq |y-z|-t|z_0-y_0|\geq|y-z|-|z_0-y_0|\geq\delta.$$
	Hence, for $t\in (0,t_0)$ it follows that
	$$\min_{z'\in\partial_{t(y_0-z_0)+x+F}(t(y_0-z_0)+K)}d_X(z')=\inf\{|y-\left(t(y_0-z_0)+z\right)\,; \, y\in\R^d\backslash X, z\in \partial_{x+F}K\}\geq \delta.$$
	By evaluating \eqref{eq: contradiction minimum principle 1} for $t\in(0,t_0)$ with $(t_0-t)|y_0-z_0|<\delta/2$, we obtain that
	$$\min_{z'\in t(y_0-z_0)+K}d_X(z')<\min_{z'\in\partial_{t(y_0-z_0)+x+F}(t(y_0-z_0)+K)}d_X(z').$$
	Since $t(y_0-z_0)+K \subseteq X \cap (t(y_0-z_0)+ x+ F)$, $d_X$ does not satisfy the minimum principle in $t(y_0-z_0)+ x+ F$, a contradiction.
\end{proof}

\end{document}
